% A research talk in Beamer with the Metropolis theme.
% Edit the metadata below, then replace the frames with your own content.
\documentclass[aspectratio=169]{beamer}

\usetheme{metropolis}
\usepackage[T1]{fontenc}
\usepackage{booktabs}
\usepackage{appendixnumberbeamer}

\title{{{title}}}
\subtitle{A short research talk}
\author{{{author}}}
\institute{Department of Economics\\University Name}
\date{\today}

\begin{document}

\begin{frame}
  \titlepage
\end{frame}

\begin{frame}{Outline}
  \tableofcontents
\end{frame}

\section{Motivation}

\begin{frame}{Why this question matters}
  Survey response rates have fallen for two decades.

  \pause

  Low response makes estimates sensitive to who answers.

  \pause

  Existing corrections rest on strong assumptions.
\end{frame}

\begin{frame}{What we do}
  \begin{itemize}
    \item<1-> Combine survey data with administrative records.
    \item<2-> Reweight the survey to match known population totals.
    \item<3-> Bound the effect under weaker assumptions.
  \end{itemize}
\end{frame}

\section{Data and design}

\begin{frame}{Two data sources}
  \begin{columns}[T, onlytextwidth]
    \column{0.48\textwidth}
    \begin{block}{Administrative records}
      Complete for every household, but coarse on income.
    \end{block}

    \column{0.48\textwidth}
    \begin{block}{Survey panel}
      Rich covariates, but response is selective.
    \end{block}
  \end{columns}
\end{frame}

\begin{frame}{Estimator}
  We estimate the coefficient by least squares:
  \begin{equation}\label{eq:ols}
    \hat{\beta} = (X^\top X)^{-1} X^\top y
  \end{equation}
  Equation~\eqref{eq:ols} holds under the classical assumptions.
\end{frame}

\begin{frame}{Three kinds of result}
  \begin{exampleblock}{Main finding}
    The effect survives reweighting.
  \end{exampleblock}

  \begin{alertblock}{Caveat}
    Bounds widen when the response model is misspecified.
  \end{alertblock}

  \begin{block}{Scope}
    The estimates describe adults in urban areas.
  \end{block}
\end{frame}

\section{Results}

\begin{frame}{Estimated effects}
  \centering
  \begin{tabular}{lcc}
    \toprule
    Specification & Estimate & Std. error \\
    \midrule
    Baseline     & 0.042 & 0.011 \\
    Reweighted   & 0.037 & 0.013 \\
    Bounded      & 0.029 & 0.015 \\
    \bottomrule
  \end{tabular}
\end{frame}

\begin{frame}{Robustness}
  The sign pattern holds across all three specifications.

  \pause

  Magnitudes shrink once the bounds are imposed.
\end{frame}

\section{Conclusion}

\begin{frame}{Takeaways}
  \begin{itemize}
    \item The main estimate is stable under reweighting.
    \item Weaker assumptions widen the interval but keep its sign.
    \item Response data should come with a sensitivity analysis.
  \end{itemize}
\end{frame}

\begin{frame}[standout]
  Questions?
\end{frame}

\appendix

\begin{frame}[noframenumbering]{Backup: alternative weights}
  Inverse probability weights give the same sign pattern.
  The standard errors are slightly larger.
\end{frame}

\end{document}
